% !TeX program = xelatex
% Four-page A4 reference sheet. Standalone; supports pdfLaTeX and XeLaTeX.
\documentclass[a4paper,10pt]{article}
\usepackage{iftex}
\ifPDFTeX
 \usepackage{cmap}\usepackage[T2A]{fontenc}\usepackage[utf8]{inputenc}
\else
 \usepackage{fontspec}
 \setmainfont{cmunrm.otf}[BoldFont=cmunbx.otf,ItalicFont=cmunti.otf,BoldItalicFont=cmunbi.otf]
\fi
\usepackage[english,russian]{babel}
\usepackage{amsmath,amssymb,mathtools,microtype}
\usepackage[left=10mm,right=10mm,top=9mm,bottom=11mm,footskip=6mm]{geometry}
\usepackage[table]{xcolor}
\usepackage[most]{tcolorbox}
\usepackage{array,tabularx,booktabs,enumitem,fancyhdr}
\usepackage[unicode,hidelinks]{hyperref}
\hypersetup{pdftitle={Алгоритмы: читлист к контрольной, 4 страницы}}
\definecolor{ink}{HTML}{354B55}
\definecolor{greenframe}{HTML}{7E9E79}
\definecolor{greenback}{HTML}{F4FBF6}
\definecolor{redframe}{HTML}{B46A55}
\definecolor{redback}{HTML}{FFF6F3}
\definecolor{purpleframe}{HTML}{9C8BB8}
\definecolor{purpleback}{HTML}{F8F6FB}
\definecolor{sandframe}{HTML}{CBBBA4}
\definecolor{sandback}{HTML}{FBF8F3}
\definecolor{blueframe}{HTML}{9CA7B3}
\definecolor{blueback}{HTML}{F4F7FA}
\pagestyle{fancy}\fancyhf{}
\renewcommand{\headrulewidth}{0pt}
\fancyfoot[L]{\fontsize{8.5}{9}\selectfont\color{ink}АСД \enspace / \enspace Справочник к контрольной}
\fancyfoot[R]{\fontsize{8.5}{9}\selectfont\color{ink}\thepage\,/\,4}
\setlength{\parindent}{0pt}
\setlength{\parskip}{2.8pt}
\setlength{\tabcolsep}{3pt}
\renewcommand{\arraystretch}{1.12}
\setlist{leftmargin=1.35em,itemsep=1pt,topsep=2pt,parsep=0pt}
\newcolumntype{Y}{>{\raggedright\arraybackslash}X}
\newtcolorbox{card}[2]{enhanced,arc=1.4mm,boxrule=.45pt,
 colback=#1back,colframe=#1frame,colbacktitle=#1frame!28!white,
 coltitle=black,title={#2},fonttitle=\bfseries\fontsize{10.5}{11.5}\selectfont,
 left=2mm,right=2mm,top=1.3mm,bottom=1.3mm,toptitle=1mm,bottomtitle=1mm,
 before skip=0pt,after skip=5pt,
 before upper={\setlength{\abovedisplayskip}{3.5pt}\setlength{\belowdisplayskip}{3.5pt}
 \setlength{\abovedisplayshortskip}{2pt}\setlength{\belowdisplayshortskip}{2pt}}}
\newcommand{\key}[1]{\textbf{#1}}
\newcommand{\Th}{\Theta}
\newcommand{\Oh}{O}
\newcommand{\Om}{\Omega}
\renewcommand{\leq}{\leqslant}
\renewcommand{\geq}{\geqslant}
\newcommand{\pagehead}[2]{%
 {\color{ink}\fontsize{17}{19}\selectfont\bfseries #1\par}
 \vspace{1mm}
 {\fontsize{9}{10}\selectfont\color{ink}#2\par}\vspace{3mm}}
\newlength{\colmax}
\setlength{\colmax}{\dimexpr\textheight-43pt\relax}
\newsavebox{\cheatcol}
\newcounter{cheatcolnum}
\newenvironment{sheetcol}[1][.486\textwidth]{%
 \begin{lrbox}{\cheatcol}\begin{minipage}[t]{#1}\vspace{0pt}
}{%
 \end{minipage}\end{lrbox}\stepcounter{cheatcolnum}%
 \typeout{CHEAT-COLUMN-\thecheatcolnum: \the\dimexpr\ht\cheatcol+\dp\cheatcol\relax; LIMIT \the\colmax}%
 \ifdim\dimexpr\ht\cheatcol+\dp\cheatcol\relax>\colmax
 \PackageError{cheat}{Column \thecheatcolnum\space exceeds planned height}{Shorten content; keep body text at least 10pt.}\fi
 \usebox{\cheatcol}}
\AtEndDocument{\clearpage\ifnum\value{page}=5\else
 \PackageError{cheat}{Expected exactly four pages}{Check the page and column layout.}\fi}
\begin{document}
\fontsize{10}{11.5}\selectfont

{\color{ink}\fontsize{9}{10}\selectfont\texttt{@hse-se-lecture-notes}\par}
\vspace{1mm}
\pagehead{1. Асимптотика: определения и переходы}
{1 --- O-символика \quad 2 --- сортировки и контейнеры \quad 3 --- циклы и корректность \quad 4 --- рекурсии}
\noindent\begin{sheetcol}
\begin{card}{green}{Что именно означает оценка\smash{\textsuperscript{*}}}
$f,g>0$ при больших $n$. Константы не зависят от $n$;
все неравенства ниже верны при $n\ge n_0$.
\[
\begin{array}{rcl}
f=O(g)&\iff&\exists C>0:\ f\le Cg,\\
f=\Omega(g)&\iff&\exists c>0:\ f\ge cg,\\
f=\Theta(g)&\iff&\exists c,C>0:\ cg\le f\le Cg.
\end{array}
\]
В каждой строке также существует порог $n_0$.
$O$ --- \key{верхняя}, $\Omega$ --- \key{нижняя},
$\Theta$ --- \key{точная по порядку} граница.

Для $g>0$:
\[
f=o(g)\iff f/g\to0,\qquad
f=\omega(g)\iff f/g\to\infty.
\]
В $o$: для \emph{каждого} $\varepsilon>0$ в итоге $f\le\varepsilon g$.
В $\omega$: для \emph{каждого} $C>0$ в итоге $f\ge Cg$.
\end{card}

\begin{card}{purple}{Равносильные переходы\smash{\textsuperscript{*}}}
\[
\begin{aligned}
f=O(g)&\iff g=\Omega(f),\\
f=\Omega(g)&\iff g=O(f),\\
f=\Theta(g)&\iff g=\Theta(f),\\
f=\Theta(g)&\iff
 \begin{cases}f=O(g),\\f=\Omega(g),\end{cases}\\
&\iff
 \begin{cases}f=O(g),\\g=O(f),\end{cases}\\
f=o(g)&\iff g=\omega(f).
\end{aligned}
\]
Для классов функций:
\[
O(f)=O(g)\iff f=\Theta(g)\iff\Omega(f)=\Omega(g).
\]
\key{Следствия:} $\Theta(g)\subseteq O(g)$ и
$\Theta(g)\subseteq\Omega(g)$.
Все пять отношений транзитивны.

\key{Не равносильно:}
$f=O(g)\ \not\Rightarrow\ g=O(f)$.
Верхняя и нижняя границы не обязаны быть точными.
\end{card}

\begin{card}{sand}{Дроби и степени}
При $f,g>0$ и постоянном $r>0$:
\[
\begin{aligned}
f=O(g)&\iff f/g=O(1)\\
&\iff 1/g=O(1/f)\\
&\iff f^r=O(g^r).
\end{aligned}
\]
При обращении дробей $O$ и $\Omega$ меняются местами:
$f=O(g)\iff 1/f=\Omega(1/g)$.
Для $\Theta$ степень и обращение также сохраняют
эквивалентность.
\end{card}

\begin{card}{blue}{Как доказать $\Theta$ по определению\smash{\textsuperscript{*}}}
Найти $c_1,c_2>0,n_0$ и показать
\[
c_1g(n)\le T(n)\le c_2g(n)\quad(n\ge n_0).
\]
Одного $T(n)\le c_2g(n)$ недостаточно:
это только $O(g)$.
\end{card}
\end{sheetcol}\hfill
\begin{sheetcol}
\begin{card}{green}{Алгебра O-символов\smash{\textsuperscript{*}}}
Для неотрицательных функций и постоянного $c>0$:
\[
\begin{aligned}
O(f)+O(g)&=O(f+g)=O(\max\{f,g\}),\\
O(f)\,O(g)&=O(fg),\\
c\,O(f)&=O(f).
\end{aligned}
\]
Те же правила верны для $\Theta$.
В конечной сумме оставляем максимальный порядок;
при произведении перемножаем порядки.
\[
f=O(g)\iff f+g=\Theta(g).
\]

$f=O(g)$ и $u=O(v)\Rightarrow fu=O(gv)$.
Сложение и умножение сохраняют также нижние оценки
для неотрицательных функций.

\key{Вычитание:} ведущие члены могут сократиться.
Из $f,g=\Theta(n)$ нельзя заключить
$f-g=\Theta(n)$.
\end{card}

\begin{card}{purple}{Сравнение через отношение}
Если существует $L=\lim f(n)/g(n)$, $g>0$:
\[
\begin{array}{c|l}
L=0 & f=o(g),\ f=O(g),\ f\notin\Omega(g)\\
0<L<\infty & f=\Theta(g)\\
L=\infty & f=\omega(g),\ f=\Omega(g),\ f\notin O(g)
\end{array}
\]
Предел может не существовать, хотя $\Theta$ верна.
Полный критерий для $f,g>0$:
\[
f=\Theta(g)\iff
0<\liminf\frac fg\le\limsup\frac fg<\infty.
\]
\end{card}

\begin{card}{sand}{Логарифмы и шкала роста\smash{\textsuperscript{*}}}
$b>1$, $a,r,k,\varepsilon$ --- постоянные:
\[
\log_b n=\frac{\ln n}{\ln b},\qquad
\log(n^r)=r\log n.
\]
\[
\log^r n=(\log n)^r\ne\log(n^r).
\]
\[
(\log n)^k=o(n^\varepsilon)\quad(k,\varepsilon>0).
\]
\[
\begin{gathered}
1\prec\log n\prec n\prec n\log n
\prec n^{3/2}\prec n^2,\\
n^k\prec n^{\log n}\prec a^n\prec n!
\quad(k>0,\ a>1).
\end{gathered}
\]
Здесь $\prec$ означает строго меньший порядок.
Полином с положительным старшим коэффициентом
имеет порядок старшей степени.
\end{card}

\begin{card}{red}{Ловушки в тестах\smash{\textsuperscript{*}}}
\key{$O\ne\Theta$.} Для $T=\Theta(n\log n)$ верны
$O(n^2)$ и $\Omega(n)$, но неверны $\Theta(n)$ и $\Omega(n^2)$.

\key{$\Omega(1)\ne O(1)$.}
$\Omega(1)$ может быть и $n$, и $2^n$.
Поэтому $O(n)+\Omega(1)$ не даёт верхней оценки $O(n)$.

\key{Случай и граница --- разные вещи.}
$O$ не означает «худший случай», а $\Omega$ --- «лучший».
Оценку применяют к уже выбранной функции времени.

\key{$f=O(g)$ --- не обычное равенство.}
Это принадлежность классу; переставлять стороны
как в числовом равенстве нельзя.
\end{card}
\end{sheetcol}
\newpage

\begingroup
\fontsize{10}{11}\selectfont
\setlength{\parskip}{1.2pt}
\renewcommand{\arraystretch}{1.0}
\pagehead{2. Сортировки и линейные структуры}
{Время --- для стандартных реализаций; сравнение ключей за $\Theta(1)$. Дополнительная память не включает входной массив.}
\begin{card}{green}{Сводная таблица сортировок}
{\fontsize{10}{11.4}\selectfont
\begin{tabularx}{\linewidth}{@{}Yllllc@{}}
\toprule
Алгоритм & Лучшее & Среднее & Худшее & Доп.\ память & Устойчив.\\\midrule
\rowcolor{greenback}
Insertion / вставки\smash{\textsuperscript{*}} & $\Theta(n)$ & $\Theta(n^2)$ & $\Theta(n^2)$ & $O(1)$ & да\\
Selection / выбором\smash{\textsuperscript{*}} & $\Theta(n^2)$ & $\Theta(n^2)$ & $\Theta(n^2)$ & $O(1)$ & нет\\
\rowcolor{greenback}
Bubble без флага\smash{\textsuperscript{*}} & $\Theta(n^2)$ & $\Theta(n^2)$ & $\Theta(n^2)$ & $O(1)$ & да\\
Bubble с флагом\smash{\textsuperscript{*}} & $\Theta(n)$ & $\Theta(n^2)$ & $\Theta(n^2)$ & $O(1)$ & да\\
\rowcolor{greenback}
Merge / слиянием\smash{\textsuperscript{*}} & $\Theta(n\log n)$ & $\Theta(n\log n)$ & $\Theta(n\log n)$ & $O(n)$ & да\\
Quick / быстрая\smash{\textsuperscript{*}} & $\Theta(n\log n)$ & $\Theta(n\log n)$ & $\Theta(n^2)$ & $O(\log n)$ / $O(n)$ & нет\\
\rowcolor{greenback}
Shell: шаги $n/2,n/4,\ldots$ & $\Theta(n\log n)$ & см.\ ниже & $\Theta(n^2)$ & $O(1)$ & нет\\
\href{https://algs4.cs.princeton.edu/cheatsheet/}{Heap / пирамидальная} & $\Theta(n\log n)^{\dagger}$ & $\Theta(n\log n)$ & $\Theta(n\log n)$ & $O(1)$ & нет\\\midrule
\href{https://ocw.mit.edu/courses/6-006-introduction-to-algorithms-fall-2011/resources/mit6_006f11_lec07_orig/}{Counting / подсчётом} & $\Theta(n+K)$ & $\Theta(n+K)$ & $\Theta(n+K)$ & $O(n+K)$ & да\\
\href{https://ocw.mit.edu/courses/6-006-introduction-to-algorithms-fall-2011/resources/mit6_006f11_lec07_orig/}{LSD Radix / поразрядная} & $\Theta(D(n+B))$ & $\Theta(D(n+B))$ & $\Theta(D(n+B))$ & $O(n+B)$ & да\\
\href{https://belohlavek.inf.upol.cz/vyuka/Cormen-bucket-sort-analysis.pdf}{Bucket / блочная} & $\Theta(n)$ & $\Theta(n)$ & $\Theta(n^2)$ & $O(n)$ & условно\\\bottomrule
\end{tabularx}
}
\key{Устойчивость:} равные ключи сохраняют исходный взаимный порядок.
Среднее для сравнительных сортировок --- по случайным перестановкам;
для quick sort также подходит случайный выбор опорного при различных ключах.
$^{\dagger}$Heap: различные ключи; с равными лучший случай зависит от реализации.
\end{card}
\setlength{\colmax}{\dimexpr\textheight-\pagetotal-10pt\relax}
\typeout{CHEAT-PAGE-2-PRECOL: \the\pagetotal; REMAINING \the\colmax}
\noindent\begin{sheetcol}[.46\textwidth]
\begin{card}{purple}{Состояние массива после прохода\smash{\textsuperscript{*}}}
\key{Insertion.} После $k$ итераций (первая обрабатывает
второй элемент) первые $k+1$ исходных элементов
отсортированы; остальной суффикс не обработан.
Число сдвигов равно числу инверсий $I$;
$T=\Theta(n+I)$, инверсия: $i<j$, $a_i>a_j$.

\key{Bubble.} При проходе слева направо максимум
уходит вправо. После $k$ проходов $k$ последних
элементов на своих местах.
Без флага раннего выхода лучший случай тоже квадратичный.

\key{Merge.} Две отсортированные части сливаются
за $\Theta(n)$; при равенстве брать из левой
части для устойчивости.
\end{card}

\begin{card}{sand}{Quick sort: опорный и разбиение\smash{\textsuperscript{*}}}
Если опорный поставлен на окончательное место $p$,
то
\[
\max_{i<p}a_i\le a_p\le\min_{i>p}a_i.
\]
Оба условия обязательны; части внутри могут быть
не отсортированы. В схеме Hoare возвращается граница,
опорный не обязан стоять на окончательном месте.

Баланс частей $\Rightarrow\Theta(n\log n)$;
размеры $0,n-1$ $\Rightarrow\Theta(n^2)$.
Стек: обычно $O(\log n)$, худший $O(n)$.
Рекурсия только в меньшую часть гарантирует
$O(\log n)$ стека.

Quick часто быстрее Merge из-за накладных расходов
и отсутствия отдельного массива, а не гарантированно
меньшего числа сравнений.
\end{card}

\begin{card}{blue}{Shell и цена сравнения}
Shell --- вставки с убывающими шагами.
Средняя оценка зависит от последовательности.
Для шагов Хиббарда / Кнута верхняя граница
$O(n^{3/2})$, для последовательности Седжвика
из семинара --- $O(n^{4/3})$.

Если сравнение стоит $O(L)$, Merge требует
$O(nL\log n)$. Для строк $L$ --- максимальная длина.
Сравнительным сортировкам в худшем случае нужно
$\Omega(n\log n)$ сравнений.
\end{card}
\end{sheetcol}\hfill
\begin{sheetcol}[.512\textwidth]
\begin{card}{green}{Массив, список, стек\smash{\textsuperscript{*}}}
Худший случай; известный узел не нужно искать.
Для массива: порядок сохраняется, без перевыделения.
\[
\begin{array}{l|ccc}
 & \text{Массив} & \text{1-связ.} & \text{2-связ.}\\ \hline
\text{Доступ по индексу} & O(1)&O(n)&O(n)\\
\text{Поиск значения} & O(n)&O(n)&O(n)\\
\text{Вставка в середину} & O(n)&O(1)^a&O(1)^a\\
\text{Удалить узел} & O(n)&O(1)^b&O(1)^a\\
\text{Удалить хвост} & O(1)&O(n)&O(1)^c
\end{array}
\]
$^a$Известно место / узел. $^b$Известен также предшественник.
$^c$Есть указатель на хвост.

Вставка в массив со сдвигом: $\Theta(n-i+1)$.
Добавление в конец при удвоении ёмкости:
$O(1)$ амортизированно, отдельный вызов $O(n)$.

В односвязном списке $k$-й от головы:
$\Theta(k)$, для фиксированного $k$ --- $\Theta(1)$;
$k$-й от конца --- $\Theta(n)$.

Стек --- LIFO; очередь --- FIFO. Операции на связной
реализации $O(1)$ (у очереди есть голова и хвост).
\end{card}

\begin{card}{sand}{Условия для линейных сортировок}
\key{Counting:} считаем частоты целых чисел.
$K$ --- число значений в диапазоне $0,\ldots,K-1$.
Время $\Theta(n+K)$; линейно при $K=O(n)$.
В таблице --- версия с выходным массивом.

\key{LSD Radix:} сортируем по цифрам от младшей
к старшей. $D$ --- число разрядов, $B$ --- основание
($B=10$ для десятичных чисел).
Линейно при постоянных $D,B$.
Проходы --- устойчивый подсчёт.

\key{Bucket:} разложить по $n$ равным интервалам,
в каждом отсортировать вставками и соединить.
Независимые равномерные данные дают $\Theta(n)$
в среднем; худшее $\Theta(n^2)$.
Для устойчивости не менять порядок равных.
\end{card}

\begin{card}{blue}{Короткие теоретические факты}
АТД --- операции и их смысл; структура данных --- реализация.

Двусвязный список: больше памяти, обход в обе стороны.
Рекурсивный вывод после обхода следующего узла:
обратный порядок; время $\Theta(n)$, стек $\Theta(n)$.

Линейный порядок: рефлексивность, антисимметричность,
транзитивность; любая пара сравнима.
\end{card}
\end{sheetcol}
\endgroup
\setlength{\colmax}{\dimexpr\textheight-43pt\relax}
\newpage

\pagehead{3. Циклы, корректность и выбор алгоритма}
{Сначала считаем реальные затраты; затем упрощаем сумму. Число уровней или итераций само по себе не равно времени.}
\noindent\begin{sheetcol}
\begin{card}{purple}{От изменения счётчика к числу шагов\smash{\textsuperscript{*}}}
Постоянные: целое $b>1$ и $d>0$.
\[
\begin{array}{l|l}
\text{Значения счётчика} & \text{Число шагов}\\ \hline
a,a+d,\ldots<M & \lceil(M-a)/d\rceil\\
1,b,b^2,\ldots<M & \lceil\log_bM\rceil\\
n,\lfloor n/b\rfloor,\ldots\ge1 & \lfloor\log_b n\rfloor+1
\end{array}
\]
Первая строка: $M>a$; вторая: $M>1$.

\key{Последовательные части:} $T=T_1+T_2$.

\key{Вложенные проходы:} $N=\sum_i N_i$.
Умножать число внешних шагов на одну внутреннюю
стоимость можно, если она одинакова на всех шагах.
\end{card}

\begin{card}{green}{Суммы, которые дают точный порядок\smash{\textsuperscript{*}}}
\[
\begin{gathered}
\sum_{i=1}^{m}i=\frac{m(m+1)}2,\\
\sum_{t=0}^{h-1}q^t=\frac{q^h-1}{q-1}\quad(q\ne1),\\
\sum_{t=0}^{\infty}q^t=\frac1{1-q}\quad(|q|<1).
\end{gathered}
\]
При $q>1$ порядок задаёт последний член;
при $0<q<1$ --- первый.
\end{card}

\begin{card}{sand}{Зависимая внутренняя работа\smash{\textsuperscript{*}}}
\key{$n$ внешних шагов; внутренние длины $1,2,4,\ldots$.}
При $h=\lceil\log_2n\rceil$:
\[
N=n\sum_{t=0}^{h-1}2^t=n(2^h-1).
\]
\[
n\le2^h<2n
\ \Rightarrow\
\tfrac12n^2\le N<2n^2\quad(n\ge2).
\]
Значит, $\Theta(n^2)$; $c_1=1/2,c_2=2,n_0=2$.

\key{Растущий счётчик; цена шага $b^t-1$.}
$h=\lceil\log_bM\rceil$, $b>1$ постоянно:
\[
N=\sum_{t=0}^{h-1}(b^t-1)
=\frac{b^h-1}{b-1}-h=\Theta(M).
\]
Если $M=cn^\alpha$, $c,\alpha>0$ постоянны,
то $T=\Theta(n^\alpha)$.
\end{card}

\begin{card}{sand}{Каскад с делением аргумента на $b$}
Для $\alpha>0$ и постоянного $b>1$:
\[
\sum_{t=0}^{\lfloor\log_b n\rfloor}(n/b^t)^\alpha
=\Theta(n^\alpha).
\]
\[
\begin{array}{c|c}
\text{Цена одного шага} & \text{Весь каскад}\\ \hline
\Theta(1) & \Theta(\log n)\\
\Theta(x^\alpha),\ \alpha>0 & \Theta(n^\alpha)\\
\Theta(\log x) & \Theta(\log^2 n)
\end{array}
\]
Цена $x,x^2,\sqrt x$ даёт порядок $n,n^2,\sqrt n$.
При цене $\log x$ суммируется арифметическая прогрессия.
\end{card}
\end{sheetcol}\hfill
\begin{sheetcol}
\begin{card}{green}{Время и число выбранных операций\smash{\textsuperscript{*}}}
\[
T(n)=\sum_j c_jN_j(n).
\]
В RAM машинная операция стоит $\Theta(1)$;
для строк и длинных чисел проверять цену отдельно.

\key{Лучшее / худшее:} минимум / максимум по входам размера $n$.
\key{Среднее:} ожидание по распределению.
\key{Амортизированное:} стоимость последовательности /
число операций; случайность не нужна.

Если при просмотре $n$ элементов запись выполняется
лишь для $r$ блоков одинаковых соседних значений,
число записей равно $r$, но полное время $\Theta(n)$,
даже при $r=1$.
\end{card}

\begin{card}{purple}{Инвариант: что проверять\smash{\textsuperscript{*}}}
\key{Инвариант $P$:} истинен перед первой итерацией
и перед каждой следующей проверкой условия.
\begin{enumerate}
\item \key{Инициализация:} предусловие $\Rightarrow P$.
\item \key{Сохранение:} $P$ и условие цикла
$\Rightarrow P$ после тела.
\item \key{Выход:} $P$ и отрицание условия
$\Rightarrow$ нужный результат.
\end{enumerate}
Завершение доказывается \key{отдельно}:
неотрицательная целая величина строго убывает.

\key{Примеры:} перед шагом $i$ сумма
$S=\sum_{j=1}^{i-1}j$; префикс до $i$ отсортирован;
в просмотренном префиксе цели нет.

Инвариантов может быть несколько.
$2=2$ --- допустимый, но бесполезный инвариант.
Внутри тела между присваиваниями $P$ может нарушаться.
\end{card}

\begin{card}{sand}{Практический выбор алгоритма\smash{\textsuperscript{*}}}
\[
T_A=an\log_b n,\quad T_B=dn,\quad a,d>0:
\]
\[
T_A<T_B\iff\log_b n<d/a\iff n<b^{d/a}.
\]
Сначала найти порог, затем сравнить с \key{допустимым диапазоном}.

Если принята модель $T=cf(n)$ и известно
$T(n_0)=t_0$, то
\[
c=\frac{t_0}{f(n_0)},\qquad
T(n)=t_0\,\frac{f(n)}{f(n_0)}.
\]
Из одного $O(f(n))$ такую точную модель получить нельзя.
\end{card}

\begin{card}{red}{Независимые повторные запуски\smash{\textsuperscript{*}}}
Успех одного запуска --- $p\in(0,1)$:
\[
P(\text{хотя бы один успех})=1-(1-p)^k.
\]
Для требуемой вероятности $q\in(0,1)$ положим
$R=\ln(1-q)/\ln(1-p)$:
\[
P\ge q:\ k_{\min}=\lceil R\rceil;\qquad
P>q:\ k_{\min}=\lfloor R\rfloor+1.
\]
Строгое «превысила» отличается от «не меньше».
При делении на $\ln(1-p)<0$ знак меняется.
Успехи должны быть независимы.
\end{card}
\end{sheetcol}
\newpage

\begingroup
\setlength{\parskip}{1.4pt}
\renewcommand{\arraystretch}{1.05}
\pagehead{4. Рекурсии: выбрать метод и получить оценку}
{Равные доли --- мастер-теорема; постоянный шаг --- «уменьшай»; разные доли --- Акра---Бацци. Всегда учитывать базу.}
\noindent\begin{sheetcol}
\begin{card}{green}{Дерево и «разделяй и властвуй»\smash{\textsuperscript{*}}}
$T(n)=aT(n/b)+f(n)$, $n=b^h$, база $\Theta(1)$:
\[
\begin{array}{l|l}
\text{Высота} & h=\log_b n\\
\text{Узлов / размер на уровне }i & a^i;\ n/b^i\\
\text{Работа уровня} & a^if(n/b^i)\\
\text{Число листьев} & a^h=n^{\log_ba}
\end{array}
\]
\[
T(n)=\sum_{i=0}^{h-1}a^if(n/b^i)+\Theta(a^h).
\]
При шаге $n-d$ высота $\Theta(n/d)$, не $\Theta(\log n)$.
Разделить $\to$ решить $\to$ объединить. Части могут
быть неравными; независимые можно распараллелить.
Эффективность не гарантирована.
\end{card}

\begin{card}{purple}{Мастер: разделяй и властвуй\smash{\textsuperscript{*}}}
$T(n)=aT(n/b)+f(n)$; постоянные $a\ge1$, $b>1$,
$f\ge0$, база $\Theta(1)$.

\key{1. Вычислить показатель рекурсии:}
\[
\boxed{p=\log_b a.}
\]
\key{2. Выделить степень $n$ в добавке:} это $k$.
Если есть множитель $\log^r n$, записать отдельно $r$.
Показатели $k,r\ge0$ постоянны.

\key{3. Сравнить $p$ и $k$.} В ячейках --- итог $T(n)$:
\[
\setlength{\arraycolsep}{3pt}
\begin{array}{c|c|c}
 & f=\Theta(n^k) & f=\Theta(n^k\log^r n)\\ \hline
p>k & \Theta(n^p) & \Theta(n^p)\\
p=k & \Theta(n^k\log n) & \Theta(n^k\log^{r+1}n)\\
p<k & \Theta(n^k) & \Theta(n^k\log^r n)
\end{array}
\]
При $p=k$ появляется \key{ещё один $\log n$}.
Если добавка задана через $O$, результаты читать как
верхние границы $O$. Для иных $f$, шага $n-d$ и разных
долей эта таблица не применяется автоматически.
\end{card}

\begin{card}{sand}{Мастер: уменьшай и властвуй\smash{\textsuperscript{*}}}
$T(n)=aT(n-d)+g(n)$, $g\ge0$; постоянные $a>0$,
целое $d>0$, база $\Theta(1)$ для всех остатков по $d$.
\key{Условия для таблицы:} $g(n)=O(F(n))$, где
$F\ge1$ не убывает. В семинаре $F(n)=n^k f(n)$.
\[
\begin{array}{c|c|c}
a & g=O(F(n)) & g=\Theta(n^k)\\ \hline
0<a<1 & O(F(n)) & \Theta(n^k)\\
a=1 & O(nF(n)) & \Theta(n^{k+1})\\
a>1 & O(a^{n/d}F(n)) & \Theta(a^{n/d})
\end{array}
\]
Средний столбец --- \key{только верхняя граница};
правый --- точная при постоянном $k\ge0$.
Например, $2T(n-1)+n$: таблица даёт $O(n2^n)$,
но точный порядок $\Theta(2^n)$.
\end{card}

\begin{card}{red}{Постоянная добавка: проверить базу\smash{\textsuperscript{*}}}
$T(n)=aT(n-d)+c$, $n=r+md$, $r$ --- база:
При $a\ne1$ положить $L=c/(1-a)$:
\[
T(r+md)=L+a^m(T(r)-L).
\]
При $a=1$: $T(r+md)=T(r)+mc$.
Если $T(r)=L$, экспонента сокращается.
Для $c<0$ таблицы выше не подходят: проверить
базу и неотрицательность решения.
\end{card}
\end{sheetcol}\hfill
\begin{sheetcol}
\begin{card}{purple}{Акра---Бацци: разные доли\smash{\textsuperscript{*}}}
\[
T(n)=\sum_{i=1}^{m}a_iT(b_i n)+g(n).
\]
$m\in\mathbb N$, $a_i>0$, $0<b_i<1$ постоянны; база $\Theta(1)$
на достаточно большом начальном промежутке.
\key{\href{https://courses.csail.mit.edu/6.046/spring04/handouts/akrabazzi.pdf}{Условие на $g$ (достаточная версия):}}
$g\ge0$ непрерывна; существуют $c_1,c_2>0$, такие что
при больших $x$ для всех $i$ и $u\in[b_i x,x]$
\[
c_1g(x)\le g(u)\le c_2g(x).
\]
Подходят $n^q\log^r n$ при постоянных $q,r\ge0$;
$2^n$ не подходит.
\key{Решение:} найти $p$ из $\sum_{i=1}^{m}a_i b_i^p=1$,
затем вычислить интеграл:
\[
\boxed{T(n)=\Theta\!\left(n^p\!\left(1+
\int_{1}^{n}\frac{g(u)}{u^{p+1}}\,du\right)\right).}
\]
\key{Пример семинара:} $2T(n/2)+\frac89T(3n/4)+n$.
$2(1/2)^2+\frac89(3/4)^2=1$, значит $p=2$;
$\int_1^n u^{-2}\,du=1-1/n$, итог $\Theta(n^2)$.
\end{card}

\begin{card}{green}{Частые рекурсии}
База положительна, $\Theta(1)$; справа от $T(n)=$:
\[
\begin{array}{l|l}
\text{Соотношение} & T(n)\\ \hline
T(n/2)+\Theta(1) & \Theta(\log n)\\
T(n/2)+\Theta(\log n) & \Theta(\log^2n)\\
T(n/2)+\Theta(n) & \Theta(n)\\
2T(n/2)+\Theta(n) & \Theta(n\log n)\\
4T(n/2)+\Theta(n^2) & \Theta(n^2\log n)\\
7T(n/2)+\Theta(n^2) & \Theta(n^{\log_2 7})
\end{array}
\]
Для Штрассена $\log_2 7\approx2.807$; на малых
матрицах меньшая степень не гарантирует выигрыш.
\end{card}

\begin{card}{sand}{Сумма размеров и сравнение\smash{\textsuperscript{*}}}
\[
T(n)=T(n/2)+T(n/4)+T(n/8)+cn,\quad c>0.
\]
$\rho=7/8<1$; работа уровня $i$ не больше $cn\rho^i$:
\[
\sum_{i\ge0}cn\rho^i=8cn;\qquad T(n)=\Theta(n).
\]
Листья дают $O(n)$; корень уже даёт $\Omega(n)$.
Длины ветвей --- от $\sim\log_8n$ до $\sim\log_2n$;
на глубоких уровнях активны \key{не все $3^i$ узлов}.
При неубывании $T$ и согласованных базах
\[
3T(n/8)+cn\le T(n)\le3T(n/2)+cn.
\]
Это даёт только $\Omega(n)$ и $O(n^{\log_2 3})$;
\key{верхняя граница неточна}. Нельзя подставлять
произвольное общее $b$ вместо разных долей.
\end{card}

\begin{card}{blue}{Подстановка = индукция\smash{\textsuperscript{*}}}
Угадать $T(n)\le Cg(n)$ $\to$ проверить базу $\to$
подставить для меньших аргументов $\to$ доказать для $n$.
Для $\Omega$ знак обратный; для $\Theta$ нужны обе границы.
Иногда помогает усиление, например $Cn-D$.
Индукция проверяет гипотезу, но не гарантирует её точность.
\end{card}
\end{sheetcol}
\endgroup
\end{document}
